\section{Counterfactual Susceptibility}
\label{sec:method}

\subsection{Single-source suppression}

Fix a prompt and an inactive target feature $i$ with positive margin
\begin{equation}
    m_i=\tau_i-z_i>0.
\end{equation}
Let $j$ be an active upstream feature with activation $a_j$. We consider suppressing its activation by a fraction $\alpha\in[0,1]$.
Our intervention convention is
\begin{equation}
    a_j \mapsto (1-\alpha)a_j,
\end{equation}
so $\alpha=0$ leaves the source unchanged and $\alpha=1$ fully suppresses it.

Let
\begin{equation}
    J_{ij}=\frac{\partial z_i}{\partial a_j}
\end{equation}
be the prompt-local response of the target preactivation to the source activation under a declared linearization convention. The first-order counterfactual target preactivation is
\begin{equation}
    \widehat z_i(\alpha)=z_i-\alpha a_jJ_{ij}.
\end{equation}

Define the movement toward the target threshold under full source suppression as
\begin{equation}
    q_{j\rightarrow i}=-a_jJ_{ij}.
\end{equation}
For an active source with positive activation, an inhibitory local relation has
$J_{ij}<0$, hence $q_{j\rightarrow i}>0$: suppressing the source raises the
predicted target preactivation. Conversely, $q_{j\rightarrow i}\leq 0$ does not
move the target toward its gate under this intervention. These statements are
prompt-local, first-order predictions; finite interventions may change active
sets and other nonlinear model states, so crossings must be verified by an
intervention sweep.
For $q_{j\rightarrow i}>0$, the predicted critical intervention is
\begin{equation}
    \widehat\alpha^\star_{j\rightarrow i}
    =\frac{m_i}{q_{j\rightarrow i}}.
\end{equation}
A full suppression reaches or exceeds the threshold in the local-linear
preactivation when $\widehat\alpha^\star_{j\rightarrow i}\leq 1$. At equality,
the prediction lands exactly at $z_i=\tau_i$; whether this activates the feature
depends on whether the loaded nonlinearity uses a strict or non-strict threshold
comparison. We therefore treat equality, and values within numerical tolerance
of it, as boundary cases rather than definite crossings.

\subsection{Dimensionless susceptibility}

We define pairwise counterfactual susceptibility as
\begin{equation}
    S_{j\rightarrow i}
    =\frac{q_{j\rightarrow i}}{m_i+\varepsilon}.
    \label{eq:susceptibility}
\end{equation}
For the unregularized score ($\varepsilon=0$), and away from numerical tolerance
boundaries, $S_{j\rightarrow i}>1$ is equivalent to
$\widehat\alpha^\star_{j\rightarrow i}<1$. With $\varepsilon>0$, however,
$S_{j\rightarrow i}>1$ requires $q_{j\rightarrow i}>m_i+\varepsilon$ and is
therefore a conservative sufficient condition for a strict full-suppression
crossing, not an exact equivalence.

The predicted target activation at intervention strength $\alpha$ is
\begin{equation}
    \widehat a_i(\alpha)
    =\phi_i\!\left(z_i-\alpha a_jJ_{ij}\right).
\end{equation}

\subsection{Behavior-weighted salience}

Let $T$ be a scalar target metric and
\begin{equation}
    g_i=\frac{\partial T}{\partial a_i}.
\end{equation}
We define a first-order behavior-weighted score
\begin{equation}
    B_{j\rightarrow i}(\alpha)
    =|g_i|\,|\widehat a_i(\alpha)-a_i|.
\end{equation}
We evaluate $S$ and $B$ separately: activation ease is not behavioral importance.

\subsection{Causal mediation}

For a verified source--target pair, let $\Delta T_j$ denote the behavior change under source suppression and let $\Delta T_{j,i0}$ denote the change when the source is suppressed while the target is clamped inactive. The signed mediated effect is
\begin{equation}
    \Delta T^{\mathrm{med}}_{j\rightarrow i}
    =\Delta T_j-\Delta T_{j,i0}.
\end{equation}
When numerically stable, we also report the mediated fraction
\begin{equation}
    \mathrm{MF}_{j\rightarrow i}
    =\frac{\Delta T^{\mathrm{med}}_{j\rightarrow i}}{\Delta T_j}.
\end{equation}

\subsection{Candidate generation}

A full dictionary search is computationally prohibitive. We therefore use a staged search: retain near-threshold inactive features, filter for positive predicted movement under selected source suppressions, and optionally prioritize features with large downstream behavioral response. Candidate recall is audited against a smaller brute-force intervention subset.
